%! Equivalent Representations of Polynomial and Rational Expressions — Unit 3, Lesson 3.5 — Group Activity (Answer Key)
\documentclass[10pt]{article}
\usepackage{precalculus-article}
\usepackage{precalculus-key}   % pulls in -boxes; do NOT also load -boxes
\newcommand{\TallMath}[1]{$\displaystyle #1\rule[-1.4em]{0pt}{3.2em}$}

\begin{document}

\pageheader{Unit 3, Lesson 3.5}{Group Activity --- Answer Key}
\namepartnerperiod

\vspace{0.1in}

\begin{scenariobox}[Three Faces of the Same Function]{plumlight}
\small
A data analyst receives four rational functions, each expressed in a different form. For each
function, the team needs to extract specific information. Sometimes the given form is ideal;
sometimes you must convert to another form first. Your job is to identify what each form reveals
and to perform the conversion when needed.

\medskip
The four functions (each given in one form) are:
\[
  f(x) = \frac{x^2 - 2x - 8}{x^2 - x - 6}
  \qquad
  g(x) = \frac{x^3 + x^2 - 2x}{x - 1}
  \qquad
  h(x) = \frac{2(x+3)(x-1)}{(x+3)(x-4)}
  \qquad
  k(x) = \frac{x^2 + 5x + 6}{x + 1}
\]
\end{scenariobox}

\vspace{0.1in}

\begin{tcolorbox}[colback=white, colframe=black!40,
    title=\textbf{Tier R --- Remediate},
    fonttitle=\bfseries, arc=2mm, left=3mm, right=3mm, top=2mm, bottom=2mm]

The numerator and denominator of $f$ and $h$ are already in factored form below.
Fill in the table using the information each form directly gives you.

\renewcommand{\arraystretch}{2.2}
\begin{tabularx}{\linewidth}{@{} c >{\centering\arraybackslash}p{3.4cm} >{\centering\arraybackslash}X >{\centering\arraybackslash}X @{}}
\textbf{Fn} & \textbf{Factored form} & \textbf{Zeros (from factored form)} & \textbf{Hole or VA at each zero?} \\
\midrule
$f$ & $\dfrac{(x-4)(x+2)}{(x-3)(x+2)}$ & \ans{$x=4,\ x=-2$} & \ans{hole at $x=-2$; VA at $x=3$} \\[6pt]
$h$ & $\dfrac{2(x+3)(x-1)}{(x+3)(x-4)}$ & \ans{$x=-3,\ x=1$} & \ans{hole at $x=-3$; VA at $x=4$} \\[6pt]
\end{tabularx}

\medskip
For each function, circle the form that best reveals end behavior:
\quad $f$: \textbf{Factored} / \ans{Standard} \qquad $h$: \textbf{Factored} / \ans{Standard}
\end{tcolorbox}

\vspace{0.1in}

\begin{tcolorbox}[colback=white, colframe=black!40,
    title=\textbf{Tier A --- Apply},
    fonttitle=\bfseries, arc=2mm, left=3mm, right=3mm, top=2mm, bottom=2mm]

\begin{enumerate}[itemsep=10pt]
  \item Convert $f(x)$ from its given form to factored form. Identify all zeros, holes, and VAs.

        \vspace{0.06in}
        Factored $f$: \ans{$\dfrac{(x-4)(x+2)}{(x-3)(x+2)}$}

        Zeros: \ans{$x=4$} \quad Hole(s): \ans{$(-2,\,6)$} \quad VA(s): \ans{$x=3$}

  \item For $g(x)$, the degree of the numerator is one more than the degree of the denominator.
        Use polynomial long division to rewrite $g(x)$ as $q(x) + \dfrac{r}{x-1}$.

        \vspace{0.06in}
        Long division work space:

        Factor numerator: $x^3+x^2-2x = x(x^2+x-2) = x(x+2)(x-1)$.

        So $g(x) = \dfrac{x(x+2)(x-1)}{x-1} = x(x+2) = x^2+2x$ for $x\ne 1$.

        \vspace{0.1in}

        $g(x) = $ \ans{$x^2 + 2x$} (remainder $= 0$, hole at $x=1$, $y=3$)

  \item For $h(x)$, determine the end behavior using the standard-form leading terms. State the
        horizontal asymptote (or explain why there is none).

        \vspace{0.06in}
        Numerator leading term: \ans{$2x^2$} \quad Denominator leading term: \ans{$x^2$}

        \vspace{0.04in}
        Horizontal asymptote: \ans{$y = 2$} \quad or \quad Slant/end behavior: \ans{HA at $y=2$; same degree}
\end{enumerate}
\end{tcolorbox}

\vspace{0.1in}

\begin{tcolorbox}[colback=white, colframe=black!40,
    title=\textbf{Tier E --- Extend},
    fonttitle=\bfseries, arc=2mm, left=3mm, right=3mm, top=2mm, bottom=2mm]

\begin{enumerate}[itemsep=10pt]
  \item Use polynomial long division on $k(x) = \dfrac{x^2 + 5x + 6}{x+1}$.
        Write $k(x) = q(x) + \dfrac{r}{x+1}$ and state the equation of the slant asymptote.

        \vspace{0.06in}
        Dividing: $x^2+5x+6 \div (x+1)$: quotient $x+4$, remainder $2$.

        $k(x) = $ \ans{$x+4+\dfrac{2}{x+1}$} \quad Slant asymptote: $y = $ \ans{$x+4$}

  \item Explain in a complete sentence how you would use the result of long division to determine
        whether $k(x)$ has a slant asymptote or a horizontal asymptote.

        \vspace{0.06in}
        \ansline{If the degree of the numerator is exactly one more than the degree of the denominator, the quotient after long division is a non-constant linear polynomial, giving a slant asymptote $y = q(x)$. If degrees are equal, the quotient is a constant, giving a horizontal asymptote.}

  \item Construct a rational function $p(x)$ that satisfies \textit{all} of the following:
        a hole at $x = 2$, a vertical asymptote at $x = -1$, zeros at $x = 0$ and $x = 5$,
        and a slant asymptote. Write $p$ in factored form and verify the degree condition for the
        slant asymptote.

        \vspace{0.06in}
        $p(x) = $ \ans{$\dfrac{x(x-5)(x-2)}{(x+1)(x-2)}$} (many valid answers)

        \vspace{0.04in}
        Degree of numerator: \ans{3} \quad Degree of denominator: \ans{2} \quad Difference: \ans{1}
\end{enumerate}
\end{tcolorbox}

\begin{teachernote}
Tier R: Students should recognize that factored form directly gives zeros and hole/VA classification without additional work.

Tier A item 2: The numerator factors cleanly, so no numeric long division is needed. Accept either the factoring approach or numeric long division.

Tier E item 3: Any rational function with numerator degree $=$ denominator degree $+ 1$, the correct zeros, hole at $x=2$, and VA at $x=-1$ is acceptable. The example shown is the minimal such construction.
\end{teachernote}

\end{document}
